\documentclass[aspectratio=169]{beamer}
\usepackage[utf8]{inputenc}
\usepackage[polish]{babel}
\usepackage[T1]{fontenc}
\usepackage{graphicx}

\usecolortheme{UMW}

% Konfiguracja nagłówka i stopki
\setbeamertemplate{navigation symbols}{}
\setbeamertemplate{footline}[frame number]

\title[Krótki Tytuł]{Tytuł Prezentacji lub Pracy Dyplomowej}
\subtitle{Podtytuł prezentacji}
\author{Imię i Nazwisko}
\institute{Uniwersytet Medyczny im. Piastów Śląskich we Wrocławiu \\ Wydział Lekarski}
\date{\today}

% Miejsce na logo uczelni
% Pobierz oficjalny plik z logo UMW i umieść w folderze jako logo_umw.png
\titlegraphic{
    %\includegraphics[height=1.5cm]{logo_umw.png}
}

\begin{document}

\begin{frame}
    \titlepage
\end{frame}

\begin{frame}{Spis treści}
    \tableofcontents
\end{frame}

\section{Wprowadzenie}

\begin{frame}{Wprowadzenie}
    \begin{itemize}
        \item Pierwszy element listy opisujący zagadnienie medyczne/naukowe.
        \item Drugi element z wyróżnieniem kluczowego pojęcia.
        \item Trzeci element podsumowujący wstęp.
    \end{itemize}
\end{frame}

\section{Metodologia i Bloki}

\begin{frame}{Formatowanie Treści}
    \begin{block}{Zwykły blok informacyjny}
        Zawiera standardowe informacje lub definicje naukowe.
    \end{block}

    \begin{alertblock}{Ważna uwaga / Wynik kluczowy}
        Blok wyróżniony złotym kolorem UMW do zwrócenia uwagi na krytyczne dane.
    \end{alertblock}

    \begin{exampleblock}{Przykład}
        Przykładowy przypadek kliniczny lub zastosowanie metody.
    \end{exampleblock}
\end{frame}

\section{Podsumowanie}

\begin{frame}{Podsumowanie}
    \begin{center}
        \Large \textbf{Dziękuję za uwagę!}
        
        \vspace{1cm}
        
        \small kontakt: imie.nazwisko@student.umw.edu.pl
    \end{center}
\end{frame}

\end{document}
